\clearpage
\begingroup
\fancyhead[L]{\small NACHLASS UND BRIEFWECHSEL}
\renewcommand{\headrulewidth}{0pt}
\par\medskip
\Needspace{8\baselineskip}
\pdfbookmark[2]{[3] Nicolai an Gauss, am Charfreitage 1814}{pallas-VIII-letter3}
\begin{center}
[3.]\par\medskip
\textsc{Nicolai} an \textsc{Gauss}, Seeberg, am Charfreitage 1814.
\end{center}
\par\noindent\hspace*{\parindent}\{.\ .\ .\ .\ .\quad Heute kann ich Ihnen wieder zwei Elemente mittheilen, und in
nächster Woche denke ich die ganze Arbeit zu endigen. Meine Besorgniss
über die zu geringe Anzahl der zum Grunde gelegten Pallasörter ist unge\SourceLeaf{588}gründet gewesen. Es können vielleicht die Gleichungen 14, 19, 20, 25\HistoricalNoteEight{Dies sind die Gleichungen mit den Argumenten:
\(4\text{\AstroSymbols ⚴}-2\text{\AstroSymbols ♄}\),
\(4\text{\AstroSymbols ⚴}-3\text{\AstroSymbols ♄}\),
\(5\text{\AstroSymbols ⚴}-3\text{\AstroSymbols ♄}\),
\(4\text{\AstroSymbols ⚴}-4\text{\AstroSymbols ♄}\).]}
beim Perihelium von der durch jene zu geringe Anzahl entstehenden Unrichtigkeit
etwas afficirt sein; allein die Coefficienten dieser Gleichungen sind
so klein, dass daraus im Allgemeinen keine wesentliche Unrichtigkeit entstehen
kann.\quad .\ .\ .\ .\ .
\par\medskip
\begingroup
\small
\newcommand{\VIIIpallas}{\text{\AstroSymbols ⚴}}
\newcommand{\VIIIsaturn}{\text{\AstroSymbols ♄}}
\setlength{\arraycolsep}{2pt}
\renewcommand{\arraystretch}{1.25}
\Needspace{12\baselineskip}
\begin{center}
Excentricitätswinkel.\par
Jährliche Änderung \(=-0\overset{\prime\prime}{,}30592\)
\end{center}
\[
\begin{array}{r@{\,}l@{\qquad}r@{\,}l}
0\overset{\prime\prime}{,}78&\sin(\VIIIpallas+191^0\ 38')&+0\overset{\prime\prime}{,}25&\sin(3\VIIIpallas-2\VIIIsaturn+320^0\ 12')\\
+1{,}52&\sin(\VIIIsaturn+187\ 49)&+0{,}60&\sin(3\VIIIsaturn+161\ 31)\\
+0{,}24&\sin(\VIIIpallas-\VIIIsaturn+202\ 49)&+0{,}76&\sin(\VIIIpallas-3\VIIIsaturn+206\ 23)\\
+0{,}14&\sin(2\VIIIpallas-\VIIIsaturn+213\ 21)&+0{,}55&\sin(2\VIIIpallas-3\VIIIsaturn+208\ 25)\\
+3{,}52&\sin(2\VIIIsaturn+206\ 12)&+0{,}30&\sin(\VIIIpallas-4\VIIIsaturn+199\ 50)\\
+0{,}16&\sin(\VIIIpallas+2\VIIIsaturn+176\ 47)&+0{,}12&\sin(2\VIIIpallas-4\VIIIsaturn+206\ 26)\\
+4{,}08&\sin(\VIIIpallas-2\VIIIsaturn+206\ 13)&+0{,}10&\sin(3\VIIIpallas-4\VIIIsaturn+205\ 41).\\
+0{,}13&\sin(2\VIIIpallas-2\VIIIsaturn+256\ 43)&&
\end{array}
\]
\Needspace{18\baselineskip}
\begin{center}
Länge des Periheliums.\par
Jährliche Änderung \(=+0\overset{\prime\prime}{,}01491\)
\end{center}
\[
\begin{array}{r@{\,}l@{\qquad}r@{\,}l}
2\overset{\prime\prime}{,}04&\sin(\VIIIpallas+345^0\ 24')&+2\overset{\prime\prime}{,}18&\sin(3\VIIIsaturn+61^0\ 21')\\
+0{,}31&\sin(2\VIIIpallas+210\ 31)&+2{,}52&\sin(\VIIIpallas-3\VIIIsaturn+296\ 44)\\
+0{,}15&\sin(3\VIIIpallas+244\ 31)&+2{,}06&\sin(2\VIIIpallas-3\VIIIsaturn+293\ 1)\\
+3{,}68&\sin(\VIIIsaturn+336\ 2)&+0{,}63&\sin(3\VIIIpallas-3\VIIIsaturn+191\ 7)\\
+0{,}23&\sin(\VIIIpallas+\VIIIsaturn+319\ 18)&+0{,}18&\sin(4\VIIIpallas-3\VIIIsaturn+232\ 6)\\
+0{,}33&\sin(\VIIIpallas-\VIIIsaturn+250\ 31)&+0{,}11&\sin(5\VIIIpallas-3\VIIIsaturn+120\ 31)\\
+0{,}52&\sin(2\VIIIpallas-\VIIIsaturn+302\ 45)&+0{,}19&\sin(4\VIIIsaturn+40\ 24)\\
+0{,}16&\sin(3\VIIIpallas-\VIIIsaturn+200\ 37)&+1{,}09&\sin(\VIIIpallas-4\VIIIsaturn+295\ 0)\\
+12{,}74&\sin(2\VIIIsaturn+109\ 32)&+0{,}42&\sin(2\VIIIpallas-4\VIIIsaturn+294\ 3)\\
+0{,}52&\sin(\VIIIpallas+2\VIIIsaturn+68\ 51)&+0{,}39&\sin(3\VIIIpallas-4\VIIIsaturn+283\ 14)\\
+13{,}77&\sin(\VIIIpallas-2\VIIIsaturn+304\ 12)&+0{,}18&\sin(4\VIIIpallas-4\VIIIsaturn+189\ 33)\\
+2{,}47&\sin(2\VIIIpallas-2\VIIIsaturn+182\ 45)&+0{,}34&\sin(\VIIIpallas-5\VIIIsaturn+303\ 19)\\
+0{,}90&\sin(3\VIIIpallas-2\VIIIsaturn+244\ 15)&+0{,}15&\sin(2\VIIIpallas-5\VIIIsaturn+292\ 20)\\
+0{,}36&\sin(4\VIIIpallas-2\VIIIsaturn+126\ 6)&+0{,}26&\sin(\VIIIpallas-6\VIIIsaturn+298\ 41).\}
\end{array}
\]
\endgroup
\clearpage
\endgroup
